% ===========================================================================
\part{The plan}
% ===========================================================================

\chapter{The defect register}
\label{ch:register}

\section{How to read it}

Two tables. The first lists \textbf{defects} --- things that are wrong in code or schema that
already exists. The second lists \textbf{gaps} --- things that are absent. The distinction matters
for sequencing: a defect gets more expensive the longer it stays, because code accumulates around
it; a gap costs nothing until you need it.

Severities mean something specific here:

\begin{description}
  \item[\severity{blocker}] Either it breaks a command you will run this week, or it will corrupt
    data, or it prevents the next piece of work from being written at all.
  \item[\severity{high}] It will cost meaningfully more to fix later than now, usually because the
    fix touches a foreign key, a wire format, or something with content already in it.
  \item[\severity{medium}] Real, worth fixing, no deadline.
  \item[\severity{low}] Housekeeping. Fix it while you are in the file.
\end{description}

Everything in both tables was confirmed by reading the code, compiling it, or running it against
the live database.

\section{Defects}

\begingroup
\small
\setlength{\LTleft}{0pt}\setlength{\LTright}{0pt}
\begin{longtable}{@{}p{0.035\textwidth}p{0.095\textwidth}L{0.375\textwidth}L{0.40\textwidth}@{}}
\caption{Defects: things that are wrong today.}\label{tab:defects}\\
\toprule
\textbf{\#} & \textbf{Sev.} & \textbf{What is wrong} & \textbf{Fix} \\
\midrule
\endfirsthead
\multicolumn{4}{@{}l}{\footnotesize\itshape Table~\ref{tab:defects} continued}\\
\toprule
\textbf{\#} & \textbf{Sev.} & \textbf{What is wrong} & \textbf{Fix} \\
\midrule
\endhead
\bottomrule
\endlastfoot

\dnum{1} & \severity{blocker} & All three foreign-key columns are declared \code{SERIAL}, so each
has its own sequence and a \code{DEFAULT nextval(...)}. An insert that omits the FK silently
attaches the child to an arbitrary parent, then later fails intermittently. Verified against the
live database. & Declare them \code{INTEGER NOT NULL REFERENCES ...}. Drop the three stray
sequences. \S\ref{sec:serialfk} \\

\dnum{38} & \severity{blocker} & \code{get\_connection\_pool()} is called \emph{inside} the
\code{HttpServer::new} factory, which runs once per worker. Verified: 12 workers on this machine
$\times$ r2d2's default \code{max\_size} of 10 = \textbf{120 connections} against a Postgres whose
\code{max\_connections} is 100. Will not fail in development. & Build the pool and \code{AppState}
above \code{HttpServer::new}; clone the state into each worker. \S\ref{sec:appstate} \\

\dnum{3} & \severity{blocker} & Two \file{down.sql} files drop the parent table before the child, so
the foreign key blocks it. \code{diesel migration revert} fails and the whole chain becomes
unrevertible. Verified. & Reverse the order: children first. \S\ref{ch:migrations} \\

\dnum{4} & \severity{blocker} & \file{create\_publications\_item/down.sql} drops
\code{publication\_itens} --- a typo. Verified: \code{ERROR: table "publication\_itens" does not
exist}. & \code{DROP TABLE IF EXISTS publication\_items;} \\

\dnum{5} & \severity{blocker} & \file{lib.rs} gates all four subject modules on \code{ssr}, so no
\code{\#[server]} function or component inside them can exist in the browser build. This now
affects real code: both \file{publications/components.rs} and the new
\file{publications/server.rs} are invisible to the wasm build. & Move the gate down to
\file{model.rs} and \file{repository.rs}. Ch.~\ref{ch:layers} \\

\dnum{6} & \severity{blocker} & \code{serde} is \code{optional = true} and enabled only by
\code{ssr}, while every model derives \code{Serialize}. Verified: the browser build fails with
\code{error[E0432]: unresolved import serde}. & Make \code{serde} unconditional.
Ch.~\ref{ch:boundary} \\

\dnum{14} & \severity{high} & \file{.env} points at database \code{my-site};
\file{docker-compose.yml} creates \code{mydatabase}. A fresh clone connects to an empty database. &
Rename both to \code{egonik\_site} (no hyphen). \S\ref{ch:deploy} \\

\dnum{7} & \severity{high} & \code{JobExperienceItem}, \code{JobInstitution},
\code{PersonalInformation} and \code{ContactInformation} have all-private fields;
\code{PortfolioItem} and \code{PublicationItem} do not. Nothing outside the module can read them. &
Add \code{pub} to every field, consistently. \S\ref{ch:dto} \\

\dnum{8} & \severity{high} & \code{Repository<T>} takes \code{\&mut self} although
\code{Pool::get} takes \code{\&self}, and its methods are synchronous. & \code{\&self}, plus
\code{web::block} at the call site. Ch.~\ref{ch:blocking} \\

\dnum{16} & \severity{high} & \code{accomplishments} and \code{responsabilities} are
\code{VARCHAR(255)} single columns holding what are logically lists. & \code{TEXT}, or a
\code{job\_highlights} child table with ordering. \\

\dnum{23} & \severity{high} & No foreign key declares \code{ON DELETE}, so the default
\code{NO ACTION} applies and deleting any parent row errors. & \code{ON DELETE CASCADE} for owned
children, \code{RESTRICT} for institutions. \\

\dnum{27} & \severity{high} & \code{contact\_informations} is one-to-one in intent and one-to-many
in the schema; nothing prevents four contact rows per person. & \code{UNIQUE (person\_id)}, or
replace with a \code{contact\_links} table. \S\ref{ch:datamodel} \\

\dnum{28} & \severity{high} & \code{job\_institutions} points the wrong way: an institution belongs
to one experience, so an employer must be duplicated per role. & Promote \code{institutions} to its
own table; \code{job\_experiences} holds the FK. \\

\dnum{24} & \severity{medium} & No index on any FK column; Postgres indexes only the referenced
side. Honest framing: at this data volume Postgres will sequentially scan regardless and should ---
the argument is schema hygiene plus the row-locking cost of \code{ON DELETE CASCADE}, not page
latency. & One \code{CREATE INDEX} per FK column. Free, so do it. \\

\dnum{19} & \severity{medium} & \file{.env} is untracked but \emph{not} in \file{.gitignore} ---
verified with \code{git check-ignore}. One \code{git add -A} commits it permanently. The exposure
today is small (the credential is \code{postgres:postgres} on a container bound to
\code{localhost}); the habit is what matters, and there is no \file{.env.example} for a fresh
clone. & Add \file{.env} to \file{.gitignore}; commit a \file{.env.example}. \\

\dnum{2} & \severity{medium} & \file{src/entrypoint/} held a \file{routes/} subtree the
compiler never saw, because \file{mod.rs} was empty. Now restructured into \file{http.rs} and
\file{application\_state.rs} --- both declared, both still empty. & Fill both, or remove them.
\S\ref{sec:appstate} \\

\dnum{9} & \severity{medium} & \code{dotenvy} and \code{anyhow} are unconditional dependencies and
are compiled into the wasm bundle. Verified with \code{cargo tree}. \code{dotenvy} reads
\file{.env} files --- in a browser. & Mark both \code{optional} and move them into the \code{ssr}
feature. \\

\dnum{11} & \severity{medium} & The portfolio, job-experience and personal-information repositories
all report \code{"Can't get publication items from db"}. Copy-paste. & Typed errors per subject; the
mistake then cannot be expressed. Ch.~\ref{ch:errors} \\

\dnum{12} & \severity{medium} & \code{get\_connection\_pool} calls \code{dotenv()}, reads
\code{env::var} and \code{.expect()}s, all inside a constructor. & Load config once in \code{main};
pass values in; return \code{Result}. \S\ref{ch:composition} \\

\dnum{13} & \severity{medium} & \file{Cargo.lock} is gitignored, though this crate produces a
binary. Builds are not reproducible. & Remove the \file{.gitignore} line; commit the lockfile. \\

\dnum{15} & \severity{medium} & \code{responsabilities} is misspelled; \code{abs} is opaque and
shadows the SQL \code{ABS()} function name. Both are about to become permanent API surface. &
\code{responsibilities}, \code{abstract\_text}. Free today. \\

\dnum{25} & \severity{medium} & No \code{CHECK} constraints. A job can end before it began; a
publication year can be 20260. & \code{CHECK (date\_to IS NULL OR date\_to >= date\_from)} and a
year range. \\

\dnum{26} & \severity{medium} & No table has \code{created\_at} / \code{updated\_at}, and the
\code{diesel\_manage\_updated\_at()} helper that the initial migration installed is unused. & Add
both columns now; they cannot be backfilled later. \\

\dnum{29} & \severity{medium} & \code{tags} stores the tag text per item, so duplicates and
near-duplicates are unavoidable and unlimited. & \code{tags} $+$ \code{portfolio\_item\_tags} join
table with a composite PK. \\

\dnum{33} & \severity{medium} & The new \code{Publications} component is \code{fn}, not
\code{pub fn}, so no route can reach it. & Add \code{pub}. \\

\dnum{10} & \severity{low} & \code{chrono} is compiled into the wasm bundle with default features,
including the parts that want a system clock and timezone database. & Keep it unconditional (dates
do cross the wire) but add \code{default-features = false, features = ["serde"]}. \\

\dnum{17} & \severity{low} & \code{println!("listening on ...")} is inside the \code{HttpServer}
factory, so it prints once per worker thread. & Move it above \code{HttpServer::new}, and use
\code{tracing::info!}. \\

\dnum{18} & \severity{low} & \code{Files::new("/assets", \&site\_root)} maps \code{/assets} onto the
entire site root, not the assets directory. Inherited from the template. &
\code{Files::new("/assets", format!("\{site\_root\}/assets"))}. \\

\dnum{20} & \severity{low} & The \code{csr} \emph{binary} target does not compile:
\code{leptos::mount\_to\_body} moved in 0.8. Verified --- \code{error[E0425]}. The lib target
compiles. & Delete the \code{csr} feature and the two dead \code{main} functions. \\

\dnum{30} & \severity{low} & \file{README.md}, \file{LICENSE} and \file{.gitignore} are still the
\code{leptos-rs/start-actix} template's; \file{Cargo.toml} has no description, authors, or
repository. & Rewrite the README as how to run \emph{this} site. \\

\dnum{32} & \severity{low} & Every migration uses \code{CREATE TABLE IF NOT EXISTS}, which turns a
schema conflict into silent divergence instead of an error. & Drop \code{IF NOT EXISTS} from
\code{CREATE}; keep \code{IF EXISTS} on \code{DROP}. \\

\dnum{21} & \severity{low} & \file{migrations/.diesel\_lock}, the \file{site} dump artefact and
\file{dumpSkip.txt} are untracked and unignored. & Add all three to \file{.gitignore}. \\

\dnum{22} & \severity{low} & Migration directory names carry a \code{-0000} suffix that Diesel does
not generate, and mix singular and plural. & Let \code{diesel migration generate} name them. \\

\dnum{31} & \severity{low} & Repositories import Diesel's internal method traits
(\code{query\_dsl::methods::\{FilterDsl, SelectDsl\}}) rather than \code{diesel::prelude::*}. &
\code{use diesel::prelude::*;} \\

\dnum{34} & \severity{low} & \code{http = \{ version = "1.3.1", optional = true \}} is enabled by no
feature and used by nothing. Verified: the only \code{http::} path in \file{src/} is
\code{actix\_web::http::}. It compiles to nothing and implies a dependency you have not got. &
Delete the line. \\

\dnum{35} & \severity{low} & \code{serde\_json} is enabled by \code{ssr} and used nowhere in
\file{src/}. Verified by \code{grep}. & Delete it until something needs it. \\

\dnum{36} & \severity{low} & \code{wasm-bindgen} and \code{console\_error\_panic\_hook} are
unconditional, so they are compiled into the \emph{native server} build as well --- the mirror image
of \dnum{9}. & Move both into the \code{hydrate} feature. \\

\dnum{37} & \severity{low} & \file{app.rs} declares both a \code{WildcardSegment} route \emph{and} a
\code{Routes fallback}. The wildcard matches first, so the fallback string \code{"Not found."} is
unreachable dead code that looks like the 404 handler. & Delete the fallback, or point it at
\code{NotFound}. \\

\end{longtable}
\endgroup

\section{Gaps}

\begingroup
\small
\setlength{\LTleft}{0pt}\setlength{\LTright}{0pt}
\begin{longtable}{@{}p{0.04\textwidth}L{0.40\textwidth}L{0.49\textwidth}@{}}
\caption{Gaps: things that do not exist yet. The order is roughly the order to build them
in.}\label{tab:gaps}\\
\toprule
\textbf{\#} & \textbf{What is missing} & \textbf{First concrete step} \\
\midrule
\endfirsthead
\multicolumn{3}{@{}l}{\footnotesize\itshape Table~\ref{tab:gaps} continued}\\
\toprule
\textbf{\#} & \textbf{What is missing} & \textbf{First concrete step} \\
\midrule
\endhead
\bottomrule
\endlastfoot

G1 & \textcolor{moss}{\textbf{Mostly done.}} \code{AppState} now holds the four repositories and
is registered with \code{app\_data}. What remains is moving its construction out of the worker
factory (\dnum{38}) and returning \code{Result} instead of panicking. & \S\ref{sec:appstate} \\

G2 & \textcolor{amber}{\textbf{Started.}} \file{publications/server.rs} has a
\code{\#[server] async fn get\_all\_publications() -> Result<(), ServerFnError>}. To become useful it
needs: a \code{Vec<PublicationDto>} return type, \code{pub}, an ungated module (\dnum{5}), and a body
that extracts \code{AppState} and calls the repository through \code{web::block}. & Ch.~\ref{ch:blocking},
\S\ref{ch:composition} \\

G3 & No page renders data. \file{app.rs} is the starter template's click counter; one real route. &
\file{src/ui/} with a \code{Shell} and one page per route. Ch.~\ref{ch:pages} \\

G4 & All seven tables are empty, and there is no seed mechanism. An empty list and a broken query
look identical. & A seed migration with your real content. Ch.~\ref{ch:seed} \\

G14 & No \code{slug} column anywhere, so detail URLs would have to be \code{/portfolio/7}. & Add
\code{slug TEXT NOT NULL UNIQUE} in the same migration as \dnum{1}. \\

G6 & No application error type and no \code{ErrorBoundary} anywhere, so a failed query renders as
nothing. (The 404 path \emph{does} work --- verified. What is missing is the 500 case.) &
\file{src/error.rs} with \code{AppError: FromServerFnError}. Ch.~\ref{ch:errors} \\

G7 & No \code{Resource}, \code{Suspense} or \code{Transition}, and no decision about SSR mode per
route. & The read pattern in Ch.~\ref{ch:loading}; pick \code{SsrMode} per route. \\

G8 & Tailwind is not set up; \file{style/main.scss} is three lines of Sass. & Four steps, no Node
required. Ch.~\ref{ch:tailwind} \\

G5 & No write operations anywhere: no \code{Insertable}, no \code{New*} struct, no
\code{insert\_into}. No authoring path, and no decision about whether there should be one. & Decide
between SQL migrations, Markdown files, or an authenticated admin. \S\ref{ch:seed} \\

G9 & No logging, no health endpoint, no graceful shutdown. One \code{println!}. &
\code{tracing} $+$ \code{tracing-actix-web}, \code{/healthz}, \code{shutdown\_timeout}.
Ch.~\ref{ch:ops} \\

G10 & No tests at all. The Playwright spec is the untouched starter example and its dependencies
are not installed. & One repository integration test; three Playwright specs.
Ch.~\ref{ch:testing} \\

G11 & No CI, no \file{rust-toolchain.toml}, no fmt/clippy gate. & The five-check workflow in
\S\ref{ch:deploy}. \\

G12 & No \file{Dockerfile}, no app service in Compose, no deployment story. & Two-stage build;
remember \code{libpq5}. Ch.~\ref{ch:deploy} \\

G13 & No SEO metadata beyond the template \code{<Title text="Welcome to Leptos"/>}: no description,
no Open Graph, no \file{robots.txt}, no sitemap. & Per-page \code{<Title>} and \code{<Meta>};
\file{robots.txt} in \file{assets/}. \\

\end{longtable}
\endgroup

\begin{pitfall}[One defect that is not in the table]
\file{src/app.rs} still sets \code{<Title text="Welcome to Leptos"/>}. It is trivial and it is also
the single most visible unfinished thing about the site --- it is what appears in a browser tab and
in search results. Worth thirty seconds before anything else in this document.
\end{pitfall}

% ===========================================================================
\chapter{The roadmap}
\label{ch:roadmap}

\section{The sequencing argument}

The temptation is to start with the visible part, because the visible part is what feels like
progress. Resist it in one specific way: \textbf{repair the foundation first}, because four of the
five blockers are things that make the next piece of work either impossible or wrong.

\begin{diagram}{Milestone dependencies. M0 unblocks everything; M1 is the one that proves the
system works end to end. M4 through M6 are genuinely independent of each other.}{fig:roadmap}
\begin{tikzpicture}[node distance=6mm]
  \node[srv,minimum width=36mm] (m0) {\textbf{M0} --- repair\\[1pt]\scriptsize schema, gates, secrets};
  \node[shared,minimum width=36mm,right=10mm of m0] (m1) {\textbf{M1} --- one slice\\[1pt]\scriptsize
    publications: DB $\rightarrow$ screen};
  \node[shared,minimum width=36mm,right=10mm of m1] (m2) {\textbf{M2} --- the rest\\[1pt]\scriptsize
    3 more slices $+$ layout};
  \node[cli,minimum width=36mm,right=10mm of m2] (m3) {\textbf{M3} --- polish\\[1pt]\scriptsize
    errors, 404, loading};

  \node[box,fill=wash,minimum width=32mm,below=11mm of m1] (m4) {\textbf{M4} --- Tailwind};
  \node[box,fill=wash,minimum width=32mm,below=11mm of m2] (m5) {\textbf{M5} --- ship\\[1pt]\scriptsize
    tracing, Docker, CI};
  \node[box,fill=wash,minimum width=32mm,below=11mm of m3] (m6) {\textbf{M6} --- optional\\[1pt]\scriptsize
    writes, SEO, tests};

  \draw[flowa] (m0) -- (m1);
  \draw[flowa] (m1) -- (m2);
  \draw[flow] (m2) -- (m3);
  \draw[flow] (m1.south) -- (m4.north);
  \draw[flow] (m1.south east) -- (m5.north west);
  \draw[flow] (m3.south) -- (m6.north);
\end{tikzpicture}
\end{diagram}

\section{M0 --- Repair the foundation}

Everything here is a defect, not a feature, and all of it is cheap \emph{today} because every table
is empty and there is one commit.

\begin{runbook}[M0, in order]
\begin{enumerate}
  \item \textbf{Secrets and hygiene first} (five minutes, and irreversible if you skip it):
    \file{.env}, \file{migrations/.diesel\_lock}, \file{site} and \file{dumpSkip.txt} into
    \file{.gitignore}; remove \file{Cargo.lock} from it and commit the lockfile. \dnum{19} \dnum{21} \dnum{13}
  \item \textbf{Rewrite the four migrations to reach the schema in
    Figure~\ref{fig:er-target}} --- one coherent target, not twelve negotiations. That single change
    covers \dnum{1} \dnum{3} \dnum{4} \dnum{15} \dnum{16} \dnum{23} \dnum{24} \dnum{25} \dnum{26}
    \dnum{27} \dnum{28} \dnum{29} \dnum{32} and G14.
  \item \textbf{Rename the database} to \code{egonik\_site} in both \file{.env} and
    \file{docker-compose.yml}, then \code{diesel database reset} and \code{diesel migration redo -a}.
    \dnum{14}
  \item \textbf{Regenerate} \file{src/schema.rs} and commit it with the migrations.
  \item \textbf{Fix the dependency table} (Table~\ref{tab:deps}): \code{serde} unconditional;
    \code{dotenvy} and \code{anyhow} \code{ssr}-only; \code{chrono} with
    \code{default-features = false}; \code{wasm-bindgen} and \code{console\_error\_panic\_hook} into
    \code{hydrate}; delete \code{http} and \code{serde\_json}; delete the \code{csr} feature and
    the two dead \code{main} functions. \dnum{6} \dnum{9} \dnum{10} \dnum{20} \dnum{34} \dnum{35} \dnum{36}
  \item \textbf{Move the module gates down} from \file{lib.rs} into each
    \file{<subject>/mod.rs}. \dnum{5}
  \item \textbf{Resolve \file{src/entrypoint/}}: put \code{AppState} in
    \file{application\_state.rs} (\S\ref{sec:appstate}) and the Actix wiring in \file{http.rs}, or
    remove both. Do not leave them empty and declared. \dnum{2}
  \item \textbf{Delete the dead \code{Routes} fallback} in \file{app.rs}. \dnum{37}
  \item \textbf{Make model fields \code{pub}}, and \code{Publications} \code{pub}. \dnum{7} \dnum{33}
  \item \textbf{Change \code{Repository} to \code{\&self}} --- or delete the trait now, since
    Ch.~\ref{ch:blocking} argues it should not survive. \dnum{8}
  \item \textbf{Fix the title.} \code{<Title text="Eduardo Gonik"/>}.
\end{enumerate}
\end{runbook}

\textbf{Done when:} \code{diesel migration redo -a} succeeds; both \code{cargo check} invocations
from Ch.~\ref{ch:twobuilds} pass; \code{git status} shows no \file{.env}.

\section{M1 --- One slice, all the way through}

\begin{keyidea}
Take \emph{one} subject --- publications, because it is the only table with no children --- and
carry it from Postgres to the rendered page. Do not start the other three until this one works.
\end{keyidea}

\begin{why}
Everything you are about to learn the hard way lives in this one slice: the gate, the DTO, the
\code{From} impl, the server function, the pool extraction, \code{web::block}, the resource, the
suspense boundary. Discovering all of that once, on the simplest table, is much cheaper than
discovering it four times in parallel --- and once it works, the other three slices are
transcription rather than design.
\end{why}

\begin{runbook}[M1, in order]
\begin{enumerate}
  \item Seed real content: three publications, via a migration. G4
  \item Move the pool and \code{AppState} construction \emph{above} \code{HttpServer::new}.
    \dnum{38} G1
  \item \file{publications/dto.rs}: \code{PublicationDto}, ungated.
  \item \file{publications/model.rs}: drop the serde derives; add \code{From<PublicationItem>}.
  \item \file{publications/repository.rs}: \code{list()} taking \code{\&self}, wrapped in
    \code{web::block}, ordered newest-first.
  \item \file{publications/api.rs}: \code{\#[server] list\_publications()}. G2
  \item \file{publications/components.rs}: grow the placeholder into the
    \code{Resource}/\code{Transition} pattern. G7
  \item Route it at \code{/publications}. G3
\end{enumerate}
\end{runbook}

\textbf{Done when:} \code{curl localhost:3000/publications} returns HTML \emph{containing the
seeded titles} --- not an empty shell that fills in later. That is the test that proves server
rendering, not just hydration, is working.

\section{M2 --- The remaining slices, and the shell}

Now repetition, plus the one genuinely new problem: the other three subjects all have child tables.

\begin{itemize}
  \item \file{src/ui/layout.rs}: \code{Shell} with header, nav, \code{<Outlet/>}, footer. Convert
    \file{app.rs} to the nested-route form. Ch.~\ref{ch:pages}
  \item \code{portfolio} with its tags, \code{job\_experience} with its institution,
    \code{personal\_information} with its contact links --- each needs an \emph{aggregate} DTO
    (\code{PortfolioItemDto} containing \code{Vec<TagDto>}), assembled in the repository with
    Diesel's \code{belonging\_to} plus \code{grouped\_by}, not with a query per row.
  \item \code{/portfolio/:slug} detail page, which is where the 404 path gets exercised for real.
\end{itemize}

\textbf{Done when:} every route in Table~\ref{tab:routes} renders seeded content, and
\code{/portfolio/nonexistent} returns HTTP 404 --- checked with \code{curl -i}, because a browser
will show you the page either way.

\section{M3 --- Errors, states, and edges}

\begin{itemize}
  \item \file{src/error.rs} with \code{AppError} implementing \code{FromServerFnError}; migrate the
    server functions off bare \code{ServerFnError}. G6
  \item \code{ErrorBoundary} on every page; one shared skeleton component with fixed heights so
    loading does not shift the layout.
  \item A real 404 page instead of the string \code{"Not found."}.
  \item Decide \code{SsrMode} per route. For a content site where crawlers matter,
    \code{SsrMode::Async} on the content routes ships complete HTML rather than a streamed shell.
\end{itemize}

\section{M4 --- Tailwind and the actual design}

Chapter~\ref{ch:tailwind} is the mechanism; this milestone is the part that is not engineering.
Do it \emph{after} M2, deliberately: styling a page whose content is real is a different and much
better experience than styling a page full of placeholder text, because the line lengths, the
heading lengths and the awkward cases are all in front of you.

\section{M5 --- Make it shippable}

\code{tracing} and \code{tracing-actix-web}; \code{/healthz}; \code{shutdown\_timeout};
embedded migrations; the two-stage \file{Dockerfile}; the app service in Compose; the five-check
CI workflow; \file{rust-toolchain.toml}. G9 G11 G12

\textbf{Done when:} \code{docker compose up} on a clean machine serves the site with content.

\section{M6 --- Whatever is actually next}

Independent, and to be picked up when each one starts to matter rather than on a schedule: the
authoring decision (G5), SEO and social previews (G13), the test suite (G10), a tags page, a blog,
an RSS feed. This is the milestone that never finishes, which is correct for a personal site.

% ===========================================================================
\chapter{Conventions and runbooks}
\label{ch:conventions}

\section{Naming}

Small decisions, made once, that stop the codebase from reading like four different projects.

\begin{table}[htbp]
\centering
\caption{Naming conventions to adopt.}
\label{tab:naming}
\small
\begin{tabular}{@{}L{0.24\textwidth}L{0.31\textwidth}L{0.36\textwidth}@{}}
\toprule
\textbf{Thing} & \textbf{Convention} & \textbf{Example} \\
\midrule
table & plural, \code{snake\_case} & \code{portfolio\_items} \\
FK column & singular parent $+$ \code{\_id} & \code{portfolio\_item\_id} \\
join table & both tables, alphabetical & \code{portfolio\_item\_tags} \\
index & \code{<table>\_<cols>\_idx} & \code{tags\_portfolio\_item\_idx} \\
Diesel model & singular, matches the row & \code{PortfolioItem} \\
DTO & model name $+$ \code{Dto} & \code{PortfolioItemDto} \\
server fn, read & \code{list\_*} / \code{get\_*} & \code{list\_portfolio}, \code{get\_portfolio\_item} \\
server fn, write & verb first & \code{create\_portfolio\_item} \\
component & \code{PascalCase}, and \code{pub} & \code{PortfolioCard} \\
migration & \code{create\_*} / \code{add\_*} / \code{fix\_*} & \code{add\_slug\_to\_portfolio\_items} \\
\bottomrule
\end{tabular}
\end{table}

Two specific points, because the current code does both ways:

\begin{itemize}
  \item \textbf{No abbreviations in schema columns.} \code{abs} saved five characters and costs
    every future reader a lookup.
  \item \textbf{Repository names should share one shape.} Today there are three:
    \code{PublicationsRepository}, \code{PortfolioItemsRepository},
    \code{JobExperienceItemRepository}. Pick \code{<Subject>Repository} and rename the other two.
\end{itemize}

\section{The language decision, which is urgent}

This one deserves its own section because it is about to become expensive, and because the codebase
has already started answering it two ways at once.

Every identifier, comment and column in the repository is in English. The one component written so
far renders Spanish:

\begin{lstlisting}[style=rs]
view! { <p> "Aca van las publicaciones" </p> }
\end{lstlisting}

Those are two \emph{different} decisions wearing the same clothes, and they should be separated
explicitly:

\begin{description}
  \item[The code's language.] English, and it already is --- with two slips
    (\code{responsabilities}, \code{publication\_itens}) that are the visible cost of drifting.
    Every dependency, every error message and every piece of documentation you will read is in
    English; matching it is the path of least friction.

  \item[The site's language.] Entirely your call, and it has no bearing on the first decision. A
    Spanish site with an English codebase is completely normal.
\end{description}

\begin{pattern}[Decide monolingual or bilingual before writing the DTOs]
This is the part that is urgent. If the site will ever be bilingual, it changes the schema and the
router, not just the strings:

\begin{itemize}
  \item the schema needs per-locale text --- either a \code{locale} column plus a composite key on
    every content table, or a \code{<table>\_translations} child table;
  \item the DTOs need to carry the resolved locale, or the server function needs a locale argument;
  \item the router needs a locale segment (\code{/es/portfolio}) or the decision to negotiate from
    the \code{Accept-Language} header;
  \item every URL slug needs a per-locale variant if you want \code{/es/portafolio}.
\end{itemize}

Retrofitting that after four subjects exist means touching every table, every DTO and every route.
Deciding it now costs one sentence. \textbf{The recommendation is monolingual} --- pick the language
your readers use and write the site in it --- because a bilingual personal site doubles the
authoring work forever, and the authoring work is the thing that actually determines whether a
personal site stays current.
\end{pattern}

\section{Runbook: adding a new content type}

The whole flow, so the next subject takes an hour rather than an afternoon of rediscovery.

\begin{runbook}[Adding a content type, nine steps]
\begin{enumerate}
  \item \code{diesel migration generate create\_talks}
  \item Write \file{up.sql}: \code{id} as \code{GENERATED ALWAYS AS IDENTITY},
    \code{slug TEXT NOT NULL UNIQUE}, the content columns, \code{position},
    \code{created\_at}/\code{updated\_at}, every FK as \code{INTEGER NOT NULL REFERENCES ... ON
    DELETE ...}, an index per FK, the \code{CHECK} constraints.
  \item Write \file{down.sql}: children first, \code{DROP TABLE IF EXISTS}.
  \item \code{diesel migration run \&\& diesel migration redo} --- \emph{never skip the redo}.
  \item \code{diesel print-schema > src/schema.rs}
  \item \file{src/talks/}: \file{mod.rs} with the gates, \file{dto.rs} (ungated),
    \file{model.rs} (\code{ssr}, with the \code{From} impl), \file{repository.rs} (\code{ssr},
    \code{\&self}, \code{web::block}), \file{api.rs} (ungated \code{\#[server]}),
    \file{components.rs} (ungated, \code{pub}).
  \item \code{pub mod talks;} in \file{lib.rs} --- \emph{ungated}.
  \item Add the route in \file{app.rs}; add a nav link in \file{ui/layout.rs}.
  \item Seed a row, then run both \code{cargo check} invocations and \code{curl} the route.
\end{enumerate}
\end{runbook}

\section{Checklist before every commit}

\begin{lstlisting}[style=sh]
cargo fmt --all
cargo clippy --no-default-features --features ssr -- -D warnings
cargo check --no-default-features --features ssr
cargo check --lib --no-default-features --features hydrate --target wasm32-unknown-unknown
# if migrations changed:
diesel migration redo -a && diesel print-schema | diff -u src/schema.rs -
git status            # is .env still untracked AND ignored?
\end{lstlisting}

\section{Five rules that summarise the whole document}

If everything else here is forgotten, these are the ones that pay for themselves.

\begin{enumerate}
  \item \textbf{Two programs, one source tree.} Every type is in the server, the browser, or both,
    and you say which. Check both builds before every commit.
  \item \textbf{Two structs per content type.} The row struct is Diesel's and stays on the server.
    The DTO is the wire's and goes everywhere. \code{From} converts.
  \item \textbf{Never block the async runtime.} Every Diesel call goes through
    \code{web::block}, with \code{pool.get()} inside the closure.
  \item \textbf{A migration is a reversible pair.} \code{diesel migration redo} after every one.
  \item \textbf{A server function is a public endpoint.} Filter in the query, not in the component;
    authorise in the body; leak nothing in the error.
\end{enumerate}
